\chapter{The Gamma and Beta Functions}

\begin{center}\itshape
The factorial is a function on the integers;\\
Euler asked which curve passes through its points, and answered with an integral.
\end{center}

In 1729 Euler wrote to Goldbach that he had found a formula giving $n!$ for every integer $n$ and making sense for
fractional $n$ as well. Two centuries later the answer to ``which curve?'' was made precise: there are infinitely many
smooth curves through the points $(n,(n-1)!)$, but only one of them has a convex logarithm. That uniqueness theorem, due
to Bohr and Mollerup (1922), is the backbone of this chapter. Following Emil Artin's little book, we use it as a machine:
to prove an identity about $\Gamma$, check that both sides satisfy the functional equation, have a convex logarithm and
agree at $1$.

The machine produces Gauss's product formula, the Beta integral, Legendre's duplication formula, and, combined with the
sine product of Chapter~3, Euler's reflection formula $\Gamma(x)\Gamma(1-x)=\frac\pi{\sin\pi x}$. At the end we keep a promise of
the Tome~0 notebooks: the volume of the unit ball in $\R^n$ is $\frac{\pi^{n/2}}{\Gamma(\frac n2+1)}$, which tends to $0$ as $n$ grows.

\begin{think}
If the factorial had a natural value at $\frac12$, what should $\bigl(\frac12\bigr)!$ be? Use $\bigl(\frac12\bigr)!^2\approx0!\cdot1!$ for a first
guess, then wait for Proposition~\ref{prop:gammahalf}.
\end{think}

\section{Euler's integral}

\begin{definition}
For $x>0$, the \emph{Gamma function} is
\[
\Gamma(x)=\int_0^\infty t^{x-1}e^{-t}\dd t.
\]
\end{definition}

The integral converges: near $0$ the integrand is at most $t^{x-1}$, integrable since $x>0$; near infinity it is
$O(t^{-2})$, since $t^{x+1}e^{-t}\to0$.

\begin{proposition}\label{prop:gammafe}
$\Gamma(1)=1$ and $\Gamma(x+1)=x\Gamma(x)$ for $x>0$. Hence $\Gamma(n+1)=n!$ for every integer $n\ge0$.
\end{proposition}
\begin{proof}
$\Gamma(1)=\int_0^\infty e^{-t}\dd t=1$. Integrating by parts on $[\varepsilon,R]$,
\[
\int_\varepsilon^Rt^xe^{-t}\dd t=\Bigl[-t^xe^{-t}\Bigr]_\varepsilon^R+x\int_\varepsilon^Rt^{x-1}e^{-t}\dd t,
\]
and the bracket tends to $0$ as $\varepsilon\to0$ and $R\to\infty$. The last statement follows by induction.
\end{proof}

\begin{proposition}\label{prop:gammasmooth}
$\Gamma$ is of class $C^\infty$ on $\left]0,\infty\right[$, with $\Gamma^{(k)}(x)=\int_0^\infty t^{x-1}(\ln t)^ke^{-t}\dd t$.
\end{proposition}
\begin{proof}
Fix $0<a<A$. For $x\in[a,A]$ and $t>0$ we have $t^{x-1}\le t^{a-1}+t^{A-1}$ (the first term dominates for $t\le1$, the second
for $t\ge1$), so
\[
\bigl|t^{x-1}(\ln t)^ke^{-t}\bigr|\le\bigl(t^{a-1}+t^{A-1}\bigr)|\ln t|^ke^{-t},
\]
which is continuous and integrable on $\left]0,\infty\right[$ for every $k$. By Theorem~\ref{thm:domination} applied $k$ times on
$J=[a,A]$, $\Gamma$ is $C^k$ there with the stated derivative. As $a$, $A$ and $k$ are arbitrary, the result follows.
\end{proof}

The functional equation extends $\Gamma$ to the left. For $x\in\left]-n,0\right[$ not an integer, set
\[
\Gamma(x)=\frac{\Gamma(x+n)}{x(x+1)\cdots(x+n-1)}.
\]
This does not depend on $n$ (replacing $n$ by $n+1$ multiplies the numerator and the denominator by $x+n$), and
$\Gamma(x+1)=x\Gamma(x)$ still holds. Near $-n$, with $n\ge0$ an integer,
\[
\Gamma(x)=\frac{\Gamma(x+n+1)}{x(x+1)\cdots(x+n)}\sim\frac{(-1)^n}{n!\,(x+n)},
\]
since the numerator tends to $\Gamma(1)=1$ and $x(x+1)\cdots(x+n-1)\to(-n)(-n+1)\cdots(-1)=(-1)^nn!$. So $\Gamma$ has a
simple pole at each of $0,-1,-2,\dots$ (Figure~\ref{fig:gamma}). It never vanishes: on $\left]0,\infty\right[$ it is the integral of a
positive function, and the formula above transports this to the left. Chapter~9 extends $\Gamma$ to the complex plane.

\begin{figure}[ht]
\centering
\fig{ch05-gamma}
\caption{$\Gamma$ on the real line. It passes through the points $(n,(n-1)!)$, has its minimum on $\left]0,\infty\right[$ at
$x_0\approx1.4616$ with $\Gamma(x_0)\approx0.8856$, and has a pole at every integer $n\le0$, where it changes sign.}
\label{fig:gamma}
\end{figure}

\section{Log-convexity and the Bohr--Mollerup theorem}

A positive function $f$ is \emph{log-convex} if $\ln f$ is convex. Log-convexity is stronger than convexity
($f=e^{\ln f}$ is then a convex function of a convex function, and $e^u$ is convex and increasing). It is stable under
products and positive powers, since $\ln(fg)=\ln f+\ln g$, and under translations $x\mapsto x+c$.

\begin{lemma}\label{lem:logconvex}
Let $I$ be an interval, $u\ge0$ and $\ell$ continuous on $I$, and $J$ an interval such that $F(x)=\int_Iu(t)e^{x\ell(t)}\dd t$
converges and is positive for every $x\in J$. Then $F$ is log-convex on $J$.
\end{lemma}
\begin{proof}
Let $x,y\in J$, $\lambda\in\left]0,1\right[$ and $p=\lambda x+(1-\lambda)y$. For $A,B\ge0$, the weighted AM--GM inequality (the concavity
of $\ln$, Tome~0, Chapter~4) gives $A^\lambda B^{1-\lambda}\le\lambda A+(1-\lambda)B$. Apply it, for each $t$, to
\[
\begin{gathered}A=\frac{u(t)e^{x\ell(t)}}{F(x)},\qq B=\frac{u(t)e^{y\ell(t)}}{F(y)},\end{gathered}
\]
for which $A^\lambda B^{1-\lambda}=\frac{u(t)e^{p\ell(t)}}{F(x)^\lambda F(y)^{1-\lambda}}$. Integrating over $I$ gives
$\frac{F(p)}{F(x)^\lambda F(y)^{1-\lambda}}\le\lambda+(1-\lambda)=1$, that is, $\ln F(p)\le\lambda\ln F(x)+(1-\lambda)\ln F(y)$.
\end{proof}

\begin{corollary}\label{cor:gammalogconvex}
$\Gamma$ is log-convex on $\left]0,\infty\right[$.
\end{corollary}
\begin{proof}
$\Gamma(x)=\int_0^\infty\frac{e^{-t}}t\,e^{x\ln t}\dd t$: apply Lemma~\ref{lem:logconvex} with $u(t)=\frac{e^{-t}}t$ and $\ell(t)=\ln t$.
\end{proof}

We use one property of a convex function $\varphi$: the slope $s(a,b)=\frac{\varphi(b)-\varphi(a)}{b-a}$ of a chord is a non-decreasing
function of each endpoint. Indeed, for $a<b<c$, write $b=\lambda a+(1-\lambda)c$ with $\lambda=\frac{c-b}{c-a}$; the convexity inequality
$\varphi(b)\le\lambda\varphi(a)+(1-\lambda)\varphi(c)$ rearranges both to $s(a,b)\le s(a,c)$ and to $s(a,c)\le s(b,c)$. In particular, for
$a<b\le c<d$, $s(a,b)\le s(b,c)\le s(c,d)$ (with the middle term omitted if $b=c$).

\begin{theorem}[Bohr--Mollerup]\label{thm:bohrmollerup}
Let $f:\left]0,\infty\right[\to\left]0,\infty\right[$ satisfy $f(1)=1$, $f(x+1)=xf(x)$ for all $x>0$, and be log-convex. Then $f=\Gamma$.
Moreover, for every $x>0$, Gauss's product formula holds:
\[
\Gamma(x)=\lim_{n\to\infty}\frac{n!\,n^x}{x(x+1)\cdots(x+n)}.
\]
\end{theorem}
\begin{proof}
Let $\varphi=\ln f$, which is convex, with $\varphi(x+1)=\varphi(x)+\ln x$; by induction $\varphi(n+1)=\ln n!$ and
\[
\varphi(x+n+1)=\varphi(x)+\ln\bigl(x(x+1)\cdots(x+n)\bigr).
\]
Let $0<x\le1$ and $n\ge1$. The three intervals $[n,n+1]$, $[n+1,n+1+x]$ and $[n+1,n+2]$ have chords of increasing slope
(the first lies to the left of the second, and the second is the third shortened on the right), so
\[
\ln n\le\frac{\varphi(n+1+x)-\varphi(n+1)}x\le\ln(n+1).
\]
Multiply by $x$, substitute the two formulas above, and let
\[
g_n(x)=\ln\frac{n!\,n^x}{x(x+1)\cdots(x+n)}.
\]
The result is $0\le\varphi(x)-g_n(x)\le x\ln\bigl(1+\frac1n\bigr)$. Hence $\varphi(x)=\lim_ng_n(x)$ for $0<x\le1$: the hypotheses determine $f$
on $\left]0,1\right]$, and then on $\left]0,\infty\right[$ through $f(x+1)=xf(x)$. Since $\Gamma$ satisfies the hypotheses
(Proposition~\ref{prop:gammafe}, Corollary~\ref{cor:gammalogconvex}), $f=\Gamma$.

Gauss's formula is proved for $0<x\le1$. A direct computation gives
\[
g_n(x+1)-g_n(x)=\ln x+\ln\frac n{x+n+1},
\]
which tends to $\ln x$. So if $g_n(x)\to\ln\Gamma(x)$, then $g_n(x+1)\to\ln\Gamma(x)+\ln x=\ln\Gamma(x+1)$, and the formula holds for all $x>0$ by induction.
\end{proof}

\begin{think}
The function $f(x)=\Gamma(x)\bigl(1+\frac3{10}\sin2\pi x\bigr)$ is positive, satisfies $f(1)=1$ and $f(x+1)=xf(x)$. Which hypothesis
of Theorem~\ref{thm:bohrmollerup} does it fail, and where? (Figure~\ref{fig:loggamma}.)
\end{think}

\begin{figure}[ht]
\centering
\fig{ch05-loggamma}
\caption{$\ln\Gamma$ (blue) is convex. The function $\ln\Gamma(x)+\ln\bigl(1+\frac3{10}\sin2\pi x\bigr)$ (orange) takes the same values at
the integers and satisfies the same functional equation, but its wiggles break convexity.}
\label{fig:loggamma}
\end{figure}

\begin{corollary}[Weierstrass's product]\label{cor:weierstrass}
For every $x>0$, with $\gamma$ Euler's constant (Tome~0, Proposition~5.5),
\[
\frac1{\Gamma(x)}=xe^{\gamma x}\prod_{k\ge1}\Bigl(1+\frac xk\Bigr)e^{-x/k},
\]
and the product converges.
\end{corollary}
\begin{proof}
With $H_n=\sum_{k\le n}\frac1k$, write $n^{-x}=e^{x(H_n-\ln n)}e^{-xH_n}$, so that
\[
\begin{aligned}\frac{x(x+1)\cdots(x+n)}{n!\,n^x}&=xe^{x(H_n-\ln n)}\eqbr\cdot\prod_{k=1}^n\Bigl(1+\frac xk\Bigr)e^{-x/k}.\end{aligned}
\]
The left side tends to $\frac1{\Gamma(x)}$ by Gauss's formula and $H_n-\ln n\to\gamma$, so the partial products converge to
$\frac{e^{-\gamma x}}{x\Gamma(x)}\ne0$.
\end{proof}

\section{Stirling's formula for the Gamma function}

Tome~0 proved $n!\sim\sqrt{2\pi n}\bigl(\frac ne\bigr)^n$ (Theorems~5.6 and~6.4). Log-convexity fills in the gaps between
integers.

\begin{theorem}\label{thm:stirlinggamma}
As $x\to\infty$ through real values, $\Gamma(x)\sim\sqrt{2\pi}\,x^{x-\frac12}e^{-x}$.
\end{theorem}
\begin{proof}
Write $x=n+t$ with $n$ an integer and $t\in[0,1[$, and let $\varphi=\ln\Gamma$. Convexity at the point $n+t=(1-t)n+t(n+1)$ gives
$\varphi(n+t)\le\varphi(n)+t\ln n$. At the point $n+1=t(n+t)+(1-t)(n+1+t)$ it gives
$\varphi(n+1)\le t\varphi(n+t)+(1-t)\varphi(n+1+t)$, and $\varphi(n+1+t)=\varphi(n+t)+\ln(n+t)$, so
\[
\varphi(n+1)\le\varphi(n+t)+(1-t)\ln(n+t).
\]
Together, with $\Gamma(n+1)=n\Gamma(n)$,
\[
\Bigl(\frac n{n+t}\Bigr)^{1-t}\le\frac{\Gamma(n+t)}{\Gamma(n)\,n^t}\le1,
\]
and the left side is at least $\frac n{n+1}$. So $\Gamma(x)=\Gamma(n)n^t(1+o(1))$, uniformly in $t$. By Tome~0,
$\Gamma(n)=\frac{n!}n\sim\sqrt{2\pi}\,n^{n-\frac12}e^{-n}$. On the other side,
\[
\begin{aligned}x^{x-\frac12}e^{-x}&=n^{n-\frac12}e^{-n}\,n^t\eqbr\cdot\Bigl(1+\frac tn\Bigr)^n e^{-t}\cdot\Bigl(1+\frac tn\Bigr)^{t-\frac12},\end{aligned}
\]
and $\bigl(1+\frac tn\bigr)^ne^{-t}\to1$ (its logarithm $n\ln(1+\frac tn)-t$ lies in $[-\frac1{2n},0]$) while the last factor lies between
$(1+\frac1n)^{-1/2}$ and $(1+\frac1n)^{1/2}$. Comparing the two expansions proves the theorem.
\end{proof}

Chapter~11 refines this into Stirling's series
$\ln\Gamma(x)=(x-\frac12)\ln x-x+\frac12\ln2\pi+\frac1{12x}-\frac1{360x^3}+\cdots$, an asymptotic series that diverges for every $x$.

\section{The Beta integral}

\begin{definition}
For $x,y>0$, the \emph{Beta function} is
\[
B(x,y)=\int_0^1t^{x-1}(1-t)^{y-1}\dd t.
\]
\end{definition}

The integral converges since $x,y>0$, and $B(x,y)=B(y,x)$ by $t\mapsto1-t$. For integers, $B(m,n)$ is the reciprocal of
an integer multiple of a binomial coefficient, and the next theorem explains why.

\begin{lemma}\label{lem:betarec}
$B(1,y)=\frac1y$ and $B(x+1,y)=\frac x{x+y}B(x,y)$ for $x,y>0$.
\end{lemma}
\begin{proof}
The first is $\int_0^1(1-t)^{y-1}\dd t=\frac1y$. For the second, integrate by parts on $[\varepsilon,1-\varepsilon]$ and let $\varepsilon\to0$; the
boundary terms vanish since $x,y>0$:
\[
\begin{aligned}B(x+1,y)&=\int_0^1t^x(1-t)^{y-1}\dd t\eqbr=\frac xy\int_0^1t^{x-1}(1-t)^y\dd t\eqbr=\frac xyB(x,y+1).\end{aligned}
\]
Writing $(1-t)^y=(1-t)^{y-1}-t(1-t)^{y-1}$ gives $B(x,y+1)=B(x,y)-B(x+1,y)$. Substituting and solving for $B(x+1,y)$ gives the
claim.
\end{proof}

\begin{theorem}[Euler]\label{thm:betagamma}
For $x,y>0$,
\[
B(x,y)=\frac{\Gamma(x)\Gamma(y)}{\Gamma(x+y)}.
\]
\end{theorem}
\begin{proof}
Fix $y>0$ and let $f(x)=\frac{\Gamma(x+y)B(x,y)}{\Gamma(y)}$ for $x>0$. Then $f>0$, and $f(1)=\frac{\Gamma(1+y)}{y\Gamma(y)}=1$ by
Lemma~\ref{lem:betarec}. With the lemma and $\Gamma(x+1+y)=(x+y)\Gamma(x+y)$,
\[
\begin{aligned}f(x+1)&=\frac{(x+y)\Gamma(x+y)}{\Gamma(y)}\cdot\frac x{x+y}B(x,y)\eqbr=xf(x).\end{aligned}
\]
Finally $x\mapsto B(x,y)=\int_0^1\frac{(1-t)^{y-1}}t\,e^{x\ln t}\dd t$ is log-convex by Lemma~\ref{lem:logconvex}, $x\mapsto\Gamma(x+y)$ is a
translate of a log-convex function, and $\Gamma(y)$ is a constant, so $f$ is log-convex. By Theorem~\ref{thm:bohrmollerup}, $f=\Gamma$.
\end{proof}

\begin{remark}
The usual proof writes $\Gamma(x)\Gamma(y)$ as a double integral over the quarter plane and changes variables $t=uv$, $s=u(1-v)$.
It needs Fubini--Tonelli and the change of variables formula for double integrals, both proved in Tome~6. Artin's proof
above needs neither: the uniqueness theorem does the work that the change of variables would do.
\end{remark}

\begin{proposition}\label{prop:gammahalf}
$\Gamma\bigl(\frac12\bigr)=\sqrt\pi$, so $\bigl(\frac12\bigr)!=\Gamma\bigl(\frac32\bigr)=\frac{\sqrt\pi}2\approx0.886$.
\end{proposition}
\begin{proof}
By Theorem~\ref{thm:betagamma}, $\Gamma(\frac12)^2=\Gamma(1)B(\frac12,\frac12)=\int_0^1\frac{\mathrm dt}{\sqrt{t(1-t)}}$. With $t=\sin^2\theta$,
$\mathrm dt=2\sin\theta\cos\theta\dd\theta$ and the integral is $\int_0^{\pi/2}2\dd\theta=\pi$. Since $\Gamma(\frac12)>0$, $\Gamma(\frac12)=\sqrt\pi$.
\end{proof}

(The substitution $t=u^2$ turns $\Gamma(\frac12)$ into $2\int_0^\infty e^{-u^2}\dd u$, so this is also the Gaussian integral of
Example~\ref{ex:gaussfeynman}, recomputed by a third method.)

\section{Reflection and duplication}

\begin{theorem}[Euler's reflection formula]\label{thm:reflection}
For every real $x$ that is not an integer,
\[
\Gamma(x)\Gamma(1-x)=\frac\pi{\sin\pi x}.
\]
\end{theorem}
\begin{proof}
First let $0<x<1$. Let $G_n(x)=\frac{n!\,n^x}{x(x+1)\cdots(x+n)}$, so that $G_n(x)\to\Gamma(x)$ by Gauss's formula, and the same at $1-x$.
In the product $G_n(x)G_n(1-x)$ the powers of $n$ combine to $n^1$. In the denominator, group the factors $k+x$ and $k-x$ for
$1\le k\le n$, which gives $\prod_k(k^2-x^2)=(n!)^2\prod_k\bigl(1-\frac{x^2}{k^2}\bigr)$; the factors $x$ and $n+1-x$ are left over. So
\[
G_n(x)G_n(1-x)=\frac n{n+1-x}\cdot\frac1{x\prod_{k=1}^n\bigl(1-\frac{x^2}{k^2}\bigr)}.
\]
As $n\to\infty$ the right side tends to $\frac\pi{\sin\pi x}$ by the sine product, Theorem~\ref{thm:sineproduct}.

For general $x$, let $R(x)=\Gamma(x)\Gamma(1-x)\sin\pi x$. By the functional equation,
$\Gamma(x+1)\Gamma(-x)=x\Gamma(x)\cdot\frac{\Gamma(1-x)}{-x}$, and $\sin\pi(x+1)=-\sin\pi x$, so $R(x+1)=R(x)$. Since $R=\pi$ on $\left]0,1\right[$,
$R=\pi$ on every interval $\left]n,n+1\right[$.
\end{proof}

At $x=\frac12$ the formula gives $\Gamma(\frac12)^2=\pi$ once more. At $x=\frac14$ it gives $\Gamma(\frac14)\Gamma(\frac34)=\pi\sqrt2$, which
Exercise~\ref{ex:lemniscate} uses to evaluate the arc length of the lemniscate.

\begin{theorem}[Legendre's duplication formula]\label{thm:duplication}
For $x>0$,
\[
\Gamma(x)\Gamma\Bigl(x+\frac12\Bigr)=2^{1-2x}\sqrt\pi\,\Gamma(2x).
\]
\end{theorem}
\begin{proof}
Let $f(x)=\frac{2^{x-1}}{\sqrt\pi}\Gamma\bigl(\frac x2\bigr)\Gamma\bigl(\frac{x+1}2\bigr)$. Then $f>0$, and $f(1)=\frac1{\sqrt\pi}\Gamma(\frac12)\Gamma(1)=1$.
Next,
\[
\begin{aligned}f(x+1)&=\frac{2^x}{\sqrt\pi}\Gamma\Bigl(\frac{x+1}2\Bigr)\Gamma\Bigl(\frac x2+1\Bigr)\eqbr=\frac{2^x}{\sqrt\pi}\Gamma\Bigl(\frac{x+1}2\Bigr)\frac x2\Gamma\Bigl(\frac x2\Bigr)\eqbr=xf(x).\end{aligned}
\]
Finally $\ln f$ is a sum of an affine function and two convex functions of the form $x\mapsto\ln\Gamma(ax+b)$ with $a>0$, so
it is convex. By Theorem~\ref{thm:bohrmollerup}, $f=\Gamma$. Replacing $x$ by $2x$ gives the formula.
\end{proof}

Exercise~\ref{ex:multiplication} extends the same proof to Gauss's multiplication formula, with $m$ factors instead of two.

\section{The volume of the unit ball}

Let $V_n$ be the volume of the unit ball $\{x\in\R^n:\ x_1^2+\cdots+x_n^2\le1\}$, with $V_0=1$. By scaling, the ball of radius $r$
has volume $r^nV_n$. Slicing the $n$-ball perpendicular to the last coordinate, the slice at height $t\in[-1,1]$ is an
$(n-1)$-ball of radius $\sqrt{1-t^2}$, so
\[
V_n=\int_{-1}^1V_{n-1}\bigl(1-t^2\bigr)^{\frac{n-1}2}\dd t.
\]
This is Cavalieri's principle: a volume is the integral of the areas of its slices. Tome~6 constructs Lebesgue measure and
proves it as a case of the Fubini--Tonelli theorem; here we take it as the definition of $V_n$, which agrees with the
familiar $V_2=\pi$ and $V_3=\frac43\pi$.

\begin{theorem}\label{thm:ballvolume}
For every $n\ge0$,
\[
V_n=\frac{\pi^{n/2}}{\Gamma\bigl(\frac n2+1\bigr)}.
\]
In particular $V_{2m}=\frac{\pi^m}{m!}$, and $V_n\to0$ as $n\to\infty$.
\end{theorem}
\begin{proof}
By symmetry and the substitution $t=\sqrt s$, for $a>-1$,
\[
\begin{aligned}\int_{-1}^1(1-t^2)^a\dd t&=\int_0^1s^{-1/2}(1-s)^a\dd s\eqbr=B\Bigl(\frac12,a+1\Bigr)\eqbr=\frac{\Gamma(\frac12)\Gamma(a+1)}{\Gamma(a+\frac32)}.\end{aligned}
\]
With $a=\frac{n-1}2$ and $\Gamma(\frac12)=\sqrt\pi$, the slicing formula becomes
\[
\frac{V_n}{V_{n-1}}=\sqrt\pi\,\frac{\Gamma(\frac{n+1}2)}{\Gamma(\frac n2+1)}.
\]
The product of these ratios for $1\le k\le n$ telescopes, since the numerator for $k$ is the denominator for $k-1$:
$V_n=\pi^{n/2}\frac{\Gamma(1)}{\Gamma(\frac n2+1)}$. Finally $\Gamma(\frac n2+1)$ grows faster than any $C^n$, by Theorem~\ref{thm:stirlinggamma}, so
$V_n\to0$.
\end{proof}

\begin{figure}[ht]
\centering
\fig{ch05-ball}
\caption{The volume $V_n$ of the unit ball in $\R^n$. It increases up to $V_5=\frac{8\pi^2}{15}\approx5.264$, then decreases to $0$:
$V_{20}\approx0.026$.}
\label{fig:ball}
\end{figure}

\begin{think}
The unit ball fits in the cube $[-1,1]^n$, of volume $2^n$, and touches every face. Yet $\frac{V_n}{2^n}\to0$ extremely fast.
Where is the volume of the cube? (Compare the distance from the centre to a corner with the radius of the ball.)
\end{think}

The maximum at $n=5$ is an accident of the unit: the ball of radius $r$ has volume $r^nV_n$, and the ratio $\frac{V_n}{V_{n-1}}$,
which decreases like $\sqrt{2\pi/n}$, falls below $1$ between $n=5$ and $n=6$. What is not an accident is the collapse:
in high dimension almost all of the cube lies in its corners, far from the inscribed ball, and almost all of the ball
lies in a thin shell near its surface (the fraction of the volume within distance $\varepsilon$ of the surface is
$1-(1-\varepsilon)^n$). This is the geometry behind the failure of grid methods and rejection sampling in high dimension met
in the notebook of Chapter~4.

\section{Exercises}

\begin{exercise}[$\star$]\label{ex:halfint}
Show that $\Gamma\bigl(n+\frac12\bigr)=\frac{(2n)!}{4^nn!}\sqrt\pi$ for every integer $n\ge0$, and compute $\Gamma\bigl(-\frac12\bigr)$ and
$\Gamma\bigl(-\frac32\bigr)$.
\end{exercise}

\begin{exercise}[$\star$]\label{ex:expp}
Show that $\int_0^\infty e^{-x^p}\dd x=\Gamma\bigl(1+\frac1p\bigr)$ for $p>0$. What is the limit as $p\to\infty$, and why is it obvious?
\end{exercise}

\begin{exercise}[$\star\star$]\label{ex:sincos}
For $a,b>-1$, show that
\[
\int_0^{\pi/2}\sin^a\theta\cos^b\theta\dd\theta=\frac12B\Bigl(\frac{a+1}2,\frac{b+1}2\Bigr),
\]
and recover the Wallis integrals $W_n=\int_0^{\pi/2}\sin^n\theta\dd\theta$ of Tome~0, Chapter~6.
\end{exercise}

\begin{exercise}[$\star\star$]\label{ex:sphere}
Define the area $S_{n-1}$ of the unit sphere in $\R^n$ as the derivative at $r=1$ of the volume $r^nV_n$ of the ball of radius $r$ (the
volume of a thin shell, divided by its thickness). Show that $S_{n-1}=nV_n=\frac{2\pi^{n/2}}{\Gamma(n/2)}$, check $S_1=2\pi$ and $S_2=4\pi$, and
find the $n$ for which $S_{n-1}$ is largest.
\end{exercise}

\begin{exercise}[$\star\star$]\label{ex:lemniscate}
The lemniscate $(x^2+y^2)^2=x^2-y^2$ has total length $2\varpi$, where $\varpi=2\int_0^1\frac{\mathrm dr}{\sqrt{1-r^4}}$ (Chapter~12). Show that
\[
\int_0^1\frac{\mathrm dr}{\sqrt{1-r^4}}=\frac{\Gamma(\frac14)^2}{4\sqrt{2\pi}}.
\]
\end{exercise}

\begin{exercise}[$\star\star\star$]\label{ex:multiplication}
Prove Gauss's multiplication formula: for every integer $m\ge1$ and $x>0$,
\[
\prod_{k=0}^{m-1}\Gamma\Bigl(x+\frac km\Bigr)=(2\pi)^{\frac{m-1}2}m^{\frac12-mx}\,\Gamma(mx).
\]
\end{exercise}

\begin{exercise}[$\star\star\star$]\label{ex:digamma}
Show that $\Gamma'(1)=-\gamma$, that is, $\int_0^\infty e^{-t}\ln t\dd t=-\gamma$, where $\gamma$ is Euler's constant.
\end{exercise}

\section{Hints}

\begin{hint}{ex:halfint}
Induction with $\Gamma(x+1)=x\Gamma(x)$, or the duplication formula at $x=n+\frac12$. For negative values, run the functional equation
backwards.
\end{hint}
\begin{hint}{ex:expp}
Substitute $t=x^p$.
\end{hint}
\begin{hint}{ex:sincos}
Substitute $t=\sin^2\theta$.
\end{hint}
\begin{hint}{ex:sphere}
Differentiate $r^nV_n$. For the maximum, look at the ratio $\frac{S_n}{S_{n-1}}$.
\end{hint}
\begin{hint}{ex:lemniscate}
Substitute $r^4=t$ to get a Beta integral, then use the reflection formula at $\frac14$.
\end{hint}
\begin{hint}{ex:multiplication}
Apply Bohr--Mollerup to $f(x)=\frac{m^{x-1}}C\prod_{k=0}^{m-1}\Gamma\bigl(\frac{x+k}m\bigr)$ with a suitable constant $C$. To find $C$, use
the reflection formula and $\prod_{k=1}^{m-1}\sin\frac{k\pi}m=\frac m{2^{m-1}}$ (Tome~0, Exercise~2.5).
\end{hint}
\begin{hint}{ex:digamma}
Take the logarithm of Weierstrass's product and differentiate term by term: you need the uniform convergence on $[\frac12,2]$
of the series of derivatives (Lemma~\ref{lem:uniform}).
\end{hint}

\section{Solutions}

\begin{solution}{ex:halfint}
\idea Each step of the functional equation multiplies by a half-integer; collecting the odd factors into a factorial gives
the closed form.
\writeup By the functional equation, $\Gamma(n+\frac12)=\frac{2n-1}2\cdot\frac{2n-3}2\cdots\frac12\,\Gamma(\frac12)=\frac{1\cdot3\cdots(2n-1)}{2^n}\sqrt\pi$, and
$1\cdot3\cdots(2n-1)=\frac{(2n)!}{2\cdot4\cdots2n}=\frac{(2n)!}{2^nn!}$. This gives the formula. Backwards, $\Gamma(-\frac12)=\frac{\Gamma(\frac12)}{-1/2}=-2\sqrt\pi$
and $\Gamma(-\frac32)=\frac{\Gamma(-\frac12)}{-3/2}=\frac43\sqrt\pi$, in accordance with the signs in Figure~\ref{fig:gamma}.
\end{solution}

\begin{solution}{ex:expp}
\idea A power substitution turns any $\int e^{-x^p}x^q\dd x$ into a Gamma integral.
\writeup With $t=x^p$, $x=t^{1/p}$ and $\mathrm dx=\frac1pt^{\frac1p-1}\dd t$, so the integral is $\frac1p\Gamma(\frac1p)=\Gamma(1+\frac1p)$. As $p\to\infty$
it tends to $\Gamma(1)=1$, since $\Gamma$ is continuous. This is obvious from the integrand: $e^{-x^p}$ tends to $1$ on $[0,1[$ and to
$0$ on $]1,\infty[$, so the area tends to that of the unit square. For $p=2$ we recover $\Gamma(\frac32)=\frac{\sqrt\pi}2$.
\end{solution}

\begin{solution}{ex:sincos}
\idea The substitution $t=\sin^2\theta$ maps $[0,\frac\pi2]$ onto $[0,1]$ and turns $\sin$ and $\cos$ into $t^{1/2}$ and $(1-t)^{1/2}$.
\writeup With $t=\sin^2\theta$, $\mathrm dt=2\sin\theta\cos\theta\dd\theta$, so
\[
\begin{aligned}\int_0^{\pi/2}\sin^a\theta\cos^b\theta\dd\theta&=\frac12\int_0^1t^{\frac{a-1}2}(1-t)^{\frac{b-1}2}\dd t\eqbr=\frac12B\Bigl(\frac{a+1}2,\frac{b+1}2\Bigr).\end{aligned}
\]
With $b=0$, $W_n=\frac12B(\frac{n+1}2,\frac12)=\frac{\sqrt\pi}2\frac{\Gamma(\frac{n+1}2)}{\Gamma(\frac n2+1)}$. For $n=2p$, Exercise~\ref{ex:halfint} gives
$W_{2p}=\frac{(2p)!}{4^p(p!)^2}\frac\pi2$, Wallis's formula. Notice that $W_n=\frac{V_n}{2V_{n-1}}$ (compare with the proof of Theorem~\ref{thm:ballvolume}): Wallis's integrals are
ratios of volumes of balls.
\end{solution}

\begin{solution}{ex:sphere}
\idea The volume of a ball is the integral of the areas of the spheres it contains, so the area is a derivative of the
volume.
\writeup $\frac{\mathrm d}{\mathrm dr}r^nV_n=nr^{n-1}V_n$, which is $nV_n$ at $r=1$. With $\frac n2\Gamma(\frac n2)=\Gamma(\frac n2+1)$ this is
$\frac{2\pi^{n/2}}{\Gamma(n/2)}$. For $n=2$, $\frac{2\pi}{\Gamma(1)}=2\pi$; for $n=3$, $\frac{2\pi^{3/2}}{\sqrt\pi/2}=4\pi$. The ratio
$\frac{S_n}{S_{n-1}}=\sqrt\pi\frac{\Gamma(\frac n2)}{\Gamma(\frac{n+1}2)}$ decreases (by log-convexity, $\frac{\Gamma(y+\frac12)}{\Gamma(y)}$ increases), and it equals
$\frac{33.07}{31.01}>1$ for $n=6$ and $\frac{32.47}{33.07}<1$ for $n=7$. So the largest sphere is $S_6$, in $\R^7$, with area
$\frac{16\pi^3}{15}\approx33.07$.
\end{solution}

\begin{solution}{ex:lemniscate}
\idea A fourth root inside a square root is a Beta integral with parameters $\frac14$ and $\frac12$; the leftover $\Gamma(\frac34)$ is
traded for $\Gamma(\frac14)$ by reflection.
\writeup With $r=t^{1/4}$, $\mathrm dr=\frac14t^{-3/4}\dd t$, the integral is $\frac14B(\frac14,\frac12)=\frac{\Gamma(\frac14)\Gamma(\frac12)}{4\Gamma(\frac34)}$. Reflection at
$x=\frac14$ gives $\Gamma(\frac34)=\frac{\pi\sqrt2}{\Gamma(1/4)}$, so the integral is
\[
\frac{\Gamma(\frac14)^2\sqrt\pi}{4\pi\sqrt2}=\frac{\Gamma(\frac14)^2}{4\sqrt{2\pi}}\approx1.31103.
\]
So $\varpi\approx2.62206$. Chapter~12 computes $\varpi$ (and so $\Gamma(\frac14)$) by Gauss's arithmetic--geometric mean.
\end{solution}

\begin{solution}{ex:multiplication}
\idea The left side, as a function of $mx$, satisfies the functional equation of $\Gamma$ up to the factor $m^x$; Bohr--Mollerup
identifies it up to a constant, and reflection pins the constant down.
\writeup Let $C=\prod_{k=1}^{m-1}\Gamma(\frac km)$ and
\[
f(x)=\frac{m^{x-1}}C\prod_{k=0}^{m-1}\Gamma\Bigl(\frac{x+k}m\Bigr).
\]
Then $f>0$, $f(1)=\frac1C\prod_{k=1}^m\Gamma(\frac km)=1$, and
$f$ is log-convex as a product of log-convex functions. Replacing $x$ by $x+1$ shifts the arguments $\frac xm,\dots,\frac{x+m-1}m$ to
$\frac{x+1}m,\dots,\frac{x+m}m$; the product gains the factor $\frac{\Gamma(\frac xm+1)}{\Gamma(\frac xm)}=\frac xm$, and $m^{x-1}$ gains a factor $m$. So
$f(x+1)=xf(x)$, and $f=\Gamma$ by Theorem~\ref{thm:bohrmollerup}.

To find $C$, pair $k$ with $m-k$ and use reflection:
\[
C^2=\prod_{k=1}^{m-1}\Gamma\Bigl(\frac km\Bigr)\Gamma\Bigl(1-\frac km\Bigr)=\frac{\pi^{m-1}}{\prod_{k=1}^{m-1}\sin\frac{k\pi}m}.
\]
Recall the proof of $\prod_{k=1}^{m-1}\sin\frac{k\pi}m=\frac m{2^{m-1}}$: with $\omega=e^{2\pi i/m}$, evaluate $\prod_{k=1}^{m-1}(X-\omega^k)=1+X+\cdots+X^{m-1}$ at
$X=1$ and use $|1-e^{i\theta}|=2\sin\frac\theta2$ for $0<\theta<2\pi$. Hence $C^2=\frac{(2\pi)^{m-1}}m$ and $C=\frac{(2\pi)^{(m-1)/2}}{\sqrt m}$, since $C>0$. Now replace $x$
by $mx$ in $\Gamma=f$:
\[
\begin{aligned}\Gamma(mx)&=\frac{m^{mx-1}\sqrt m}{(2\pi)^{(m-1)/2}}\prod_{k=0}^{m-1}\Gamma\Bigl(x+\frac km\Bigr),\end{aligned}
\]
which rearranges to the formula. For $m=2$ it is Legendre's duplication formula.
\end{solution}

\begin{solution}{ex:digamma}
\idea The logarithm turns Weierstrass's product into a series, whose derivative at $1$ telescopes.
\writeup Taking logarithms in Corollary~\ref{cor:weierstrass},
\[
\begin{gathered}-\ln\Gamma(x)=\ln x+\gamma x+\sum_{k\ge1}u_k(x),\qq u_k(x)=\ln\Bigl(1+\frac xk\Bigr)-\frac xk.\end{gathered}
\]
On $[\frac12,2]$, $u_k'(x)=\frac1{k+x}-\frac1k=-\frac x{k(k+x)}$ satisfies $|u_k'(x)|\le\frac2{k^2}$, so the series of derivatives converges uniformly
(Weierstrass test), and the series itself converges at $x=1$. By Lemma~\ref{lem:uniform} the sum is differentiable with the
termwise derivative, so
\[
\begin{aligned}-\frac{\Gamma'(x)}{\Gamma(x)}&=\frac1x+\gamma+\sum_{k\ge1}\Bigl(\frac1{k+x}-\frac1k\Bigr).\end{aligned}
\]
At $x=1$ the series telescopes to $\sum_k(\frac1{k+1}-\frac1k)=-1$, so $-\Gamma'(1)=1+\gamma-1=\gamma$. By
Proposition~\ref{prop:gammasmooth}, $\Gamma'(1)=\int_0^\infty e^{-t}\ln t\dd t$.
\end{solution}

\section*{Sources and further reading}

\begin{itemize}
\item E.~Artin, \emph{The Gamma Function}, translated by M.~Butler, Holt, Rinehart and Winston, 1964 (German original 1931).
Thirty-nine pages; the proofs of Theorems~\ref{thm:betagamma} and~\ref{thm:duplication} follow it.
\item H.~Bohr and J.~Mollerup, \emph{L\ae{}rebog i matematisk Analyse}, vol.~III, Copenhagen, 1922, where
Theorem~\ref{thm:bohrmollerup} first appeared.
\item P.~J.~Davis, \emph{Leonhard Euler's integral: a historical profile of the Gamma function}, American Mathematical Monthly
66 (1959), 849--869.
\item G.~Andrews, R.~Askey and R.~Roy, \emph{Special Functions}, Cambridge University Press, 1999, Chapter~1.
\item R.~Remmert, \emph{Wielandt's theorem about the $\Gamma$-function}, American Mathematical Monthly 103 (1996), 214--220. The
complex analogue of Bohr--Mollerup, for Chapter~9.
\item C.~Lanczos, \emph{A precision approximation of the gamma function}, SIAM Journal on Numerical Analysis, Series~B, 1
(1964), 86--96. The method of the notebook.
\end{itemize}
